\chapter{AnchorEvolve：前瞻性断言分支}
\label{ch:anchor}

AnchorEvolve是FinPlanRAG的前瞻性断言分支。本章先给出状态表示、证据来源和总体结构，再给出完整算法，包括证据构造、锚点抽取、五类资格判定、时间演化、类型库扩展、信念更新、全动作读出、数值桥和披露日期断言。

% ==================== 第4章 ====================

\section{设计原则}

FinPlanRAG的设计遵循六条原则，其核心模块AnchorEvolve承担证据状态化、资格判定、时间演化和信念更新。这些原则不是附加的功能列表，而是对第\ref{ch:framework}章所述研究缺口的直接回应。

第一，契约先行。系统在检索证据之前先确定公司、目标期间、指标、结果空间和截止时点。契约把“回答什么”写成显式对象，使不同方法可以在同一口径下比较，也使用户可以在运行前检查阈值是否合理。契约不是系统的内部细节，而是审计报告的第一部分。

第二，来源落地。每条锚点必须引用一个具体的证据来源，并给出逐字引文。系统不接受没有来源的判断，也不接受只在理由中描述证据而不给出引文的锚点。来源落地使锚点可以被逐条复核，也使错误的锚点可以被定位到具体的网页或段落。

第三，资格先于权重。锚点在进入信念更新之前先经过资格判定。空泛、过期、对抗三类锚点不进入更新；硬锚点和软锚点分别处理。这个顺序很重要：如果先把所有锚点加权再过滤，被拒绝的锚点仍然会影响后验；如果先过滤再加权，每个被采纳的锚点都有明确的资格理由。

第四，揭示后演化。锚点类型的可靠性只在结果公开揭示之后更新，且只使用揭示日早于当前截止时点的结果。这个约束使可靠性演化在回测和部署中具有相同的含义，避免把未来结果泄漏到时点判断。

第五，损失感知。输出不是固定的点答案，而是在点答案、非单元素裁决集和弃权三个动作上最小化期望损失。损失函数把过度乐观、漏报、集合宽度和弃权成本写成显式参数，使不同风险偏好可以在同一框架下表达。

第六，审计优先。系统把证据、锚点、资格、可靠性、信念轨迹、动作损失和日期估计写入结构化文件。审计不是事后附加的日志，而是系统输出的一部分。没有审计文件的判断不能被本文视为可复核的判断。

\section{三层锚点状态}

FinPlanRAG的核心状态是AnchorEvolve的三层锚点。三层分别是稳定关系层（L0）、行业模式层（L1）和实例层（L2）。这里的锚点不是知识图谱中的实体或关系，也不是行为金融中的anchoring bias；L0、L1和L2分别对应关系约束、行业模式绑定和带来源的证据实例，锚定的是证据对结果的约束关系，而不是一个实体节点。AnE用“anchor”指训练数据的真值锚定，本文的锚点指证据状态中的结构化对象，二者机制不同\citep{wang2026ane}。

L0是稳定关系本体。它定义一组可复用的关系类型，包括界、先后、前置时间、持续时间、守恒、比率、单位换算和作用域。L0不随行业变化，为锚点提供语义骨架。例如，“当前季度值落在指引区间内”是一个界关系，“新产品在指引季度之后投产”是一个先后关系，“产能爬坡需要两个季度”是一个前置时间关系，“分部收入之和等于总收入”是一个守恒关系，“毛利率等于毛利润除以营业收入”是一个比率关系。L0的作用是限制锚点可以表达的关系种类，使不同行业的锚点共享同一套语义。

L1是行业模式。每个锚点类型是一个四元组
\begin{equation}
\tau=(\text{名称},\ \text{L0关系},\ \text{字段},\ \text{生效轮次}),
\label{eq:l1}
\end{equation}
其中字段给出该类型需要从证据中读取的信息，生效轮次给出该类型可以在哪一轮断言中使用。收入指引的初始模式包括指引宽度、当前季度增长、公司历史、行业历史、文本信号、产能与需求、分析师预测、披露历史和监管规则。L1把L0关系绑定到具体字段和生效规则。例如，“指引宽度”绑定到界关系，字段是上限与下限之差除以中点；“当前季度增长”绑定到比率关系，字段是当前季度值相对指引中点的比例；“披露历史”绑定到先后关系，字段是公司过去若干年年报的公告日期。

L2是锚点实例。每条锚点是一个元组
\begin{equation}
a=(id,\tau,v,u,s,src,d,r,\text{hard},\hat h,\text{note}),
\label{eq:l2}
\end{equation}
其中$id$是锚点编号，$\tau$是L1类型，$v$是取值或区间，$u$是单位，$s$是作用域，$src$是来源编号，$d$是来源日期，$r$是可靠性，$\text{hard}$是硬锚点标记，$\hat h$是隐含结果标签，$\text{note}$是理由。L2是LLM抽取器的直接输出，也是资格判定和信念更新的输入。每条L2锚点都记录来源和时间，因此可以被逐条检查。

三层的关系可以概括为：L0定义可以说哪些关系，L1定义某类断言需要哪些字段，L2记录某条证据的具体取值。三层不是三个独立训练的模块，而是同一状态的不同抽象层次。L1类型用于可靠性演化，L0关系用于语义分类，L2实例进入信念更新。这样的分层使系统可以在不改变L0的情况下增加新的L1类型，也可以在不改变L1的情况下替换L2抽取器。

\section{锚点类型}

表\ref{tab:anchortypes}给出本文使用的主要锚点类型。类型的选择遵循两个标准：一是该类型在金融文本中有稳定的语言表达，二是该类型能够对应到结果空间中的一个或多个标签。类型不是越多越好；过多的类型会增加抽取器的负担，也会增加可靠性后验的稀疏性。

\begin{table}[H]
\centering
\zihao{5}
\caption{FinPlanRAG的主要锚点类型。}
\label{tab:anchortypes}
\begin{tabular}{p{2.6cm}p{2.4cm}p{7.2cm}}
\toprule
类型 & L0关系 & 作用 \\
\midrule
指引宽度 & 界 & 指引区间越宽，结果落在区间内的先验越高；区间本身不是结果落在区间内的证据。\\
当前季度增长 & 比率 & 当前季度实际值相对指引中点的比例，用于判断全年结果的起点。\\
公司历史 & 先后 & 公司过去已揭示结果相对指引的偏差，用于估计公司层面的兑现倾向。\\
行业历史 & 先后 & 同行业其他公司已揭示结果的分布，用于估计行业层面的共同冲击。\\
文本信号 & 语义 & 管理层评论、需求描述、订单和产能信息，用于补充结构化字段。\\
产能与需求 & 前置时间 & 产能投放、利用率和订单覆盖，用于判断收入增长的可持续性。\\
分析师预测 & 界 & 卖方机构对目标期间的预测，作为软锚点而不是事实。\\
披露历史 & 先后 & 公司过去年度报告的公告日期，用于披露日期断言。\\
监管规则 & 界 & 交易所的披露截止日，作为上界而不是具体日期证据。\\
\bottomrule
\end{tabular}
\end{table}

锚点类型与结果标签的关系不是一一对应。一条“当前季度增长”锚点可能支持“高于”标签，也可能只支持“不低于”某个下界；一条“指引宽度”锚点通常不直接支持任何标签，而是影响区间内的先验。资格判定和信念更新需要分别处理这些情况。锚点的$\hat h$字段记录抽取器认为该锚点隐含的标签；如果锚点不提供标签信息，$\hat h$为空，锚点被判为空泛。

\section{证据来源与作用域}

FinPlanRAG使用三类证据来源。

第一类是公司申报文件和公告。美国公司使用8-K、10-Q、10-K和XBRL数据；中国A股公司使用巨潮资讯网的公告记录、年报、季报和临时公告。申报文件提供结构化事实和原文段落，公告记录提供披露日期和文件链接。官方来源的优先级高于媒体转载，因为官方来源的日期和文本更可靠。

第二类是财经媒体和研究报告。媒体文章可以提供管理层评论、行业需求、产能和价格信息，但可能存在转载、改编和日期错误。系统要求媒体来源给出逐字引文，并对无日期来源降权。研究报告中分析师预测作为软锚点，不作为事实。

第三类是交易所规则和监管文件。这类来源提供披露截止日、披露格式和合规要求。监管规则通常只给出上界，不能直接支持某个具体日期窗口；系统把泛泛的监管截止日判为空泛锚点，除非来源明确给出公司的预约披露日期。

锚点的作用域分为公司、行业、市场和监管四类。作用域的作用是防止跨作用域误用。例如，同行业其他公司的营收指引结果可以作为行业锚点，但不能直接当作目标公司的结果；交易所的披露截止日可以作为监管锚点，但不能直接推断公司会在截止日当天披露。资格判定检查作用域是否与断言匹配，不匹配的锚点被判为对抗锚点。

\section{与既有框架的区别}

FinPlanRAG与断言验证、RAG和选择性预测的区别可以概括为三点。

第一，状态不同。断言验证的输出是一个离散裁决，RAG的输出是一个生成答案，选择性预测的输出是一个标签或弃权。FinPlanRAG的输出是一个可演化的锚点状态，包含证据、资格、可靠性和信念。这个状态可以被逐条检查，也可以在结果揭示后更新。

第二，时间不同。断言验证和RAG通常在一次调用中完成，不区分证据的可用时间和结果的揭示时间。FinPlanRAG把截止时点写进契约，把揭示日写进可靠性更新规则。这个区分使系统可以在回测中复现部署时的信息条件。

第三，决策不同。断言验证和RAG通常报告准确率或F1，选择性预测报告覆盖率或弃权率。FinPlanRAG在同一个损失下比较点答案、裁决集和弃权，报告决策损失、覆盖率和集合大小。这个口径使系统可以在不确定性高时选择集合或弃权，而不是被迫输出一个点答案。

% ==================== 第5章 ====================

本章给出FinPlanRAG中AnchorEvolve模块的完整算法设计。本章第1节到第7节分别介绍证据构造、锚点抽取、资格判定、时间演化、信念更新、全动作读出和披露日期断言，最后给出算法伪代码、复杂度和审计格式。

\section{证据构造}

证据构造的目标是为每条断言提供一组带日期、带来源编号、可引用的材料。构造过程分为查询生成、检索、去重、日期过滤和编号五步。

查询生成从契约出发。给定公司$e$、目标期间$p$和指标$m$，系统生成四类查询：公司加期间加指标，公司加最近季度加指标，公司加上一年度加指标，以及公司加“预测”“研报”“公告”等限定词。对于披露日期断言，系统额外生成公司加“年度报告披露日期”、公司加“往年披露时间”、公司加“预约披露”等查询。查询数量由配置限制，默认每条断言最多四条。查询中不放入断言的目标答案，避免把判断泄漏到检索。

检索使用多个网页搜索提供者，并按轮换顺序调用。某个提供者返回空结果或安全验证页时，系统自动切换到下一个提供者。这个设计是为了避免单一搜索引擎的限流让整次运行静默失败。检索结果包含标题、链接、摘要、发布日期和提供者。对于披露日期断言，系统还查询巨潮资讯网的公告接口，直接获取公司历年年报的公告日期和PDF链接。官方公告记录作为高优先级证据，排在网页结果之前。

去重按链接和标题进行。同一篇文章的不同转载只保留一条；同一链接的不同摘要只保留一条。日期过滤检查证据的发布日期是否晚于截止时点。晚于截止时点的证据被写入剔除列表，不进入锚点抽取。发布日期未知的证据保留，但标记为无日期，并在资格判定中降权。

编号使用$W1,W2,\dots$，按证据得分排序。得分由日期是否已知、摘要是否包含数字、标题或摘要是否包含公司名、文本长度和查询顺序决定。官方公告记录得分最高。编号的作用是让锚点可以引用具体来源，也让审计报告可以把锚点映射回证据。

对于申报文件，系统在证据构造阶段还做段落切分。每个段落带稳定编号$P1,P2,\dots$，结构化证据带$E1,E2,\dots$。抽取器可以引用$P\#$或$E\#$，资格判定检查引用的段落是否在时点摘录中。这个设计使锚点可以定位到段落，而不是只定位到文件。

\section{锚点抽取}

锚点抽取是一个LLM调用。输入包括断言契约、结果标签及其定义、参考口径、时点证据、锚点类型和关系类型。输出是一个JSON对象，包含锚点列表、标签分数、日期估计（仅日期断言）、置信度和缺失说明。

提示词要求锚点遵守以下规则。第一，每条锚点必须引用一个证据编号，不能引用不存在的来源。第二，每条锚点必须给出逐字引文，引文必须在来源文本中出现。第三，锚点的发布日期不得晚于截止时点；来源日期未知时，日期设为空，可靠性不得高于0.55。第四，可靠性是“隐含标签正确的校准概率”，不是推理置信度。第五，硬锚点必须是显式约束，能够排除至少一个标签；管理层评论和分析师预测通常是软锚点。第六，证据不足时在缺失字段中说明，不编造事实。第七，标签分数必须使用契约中的标签，且和为1。

抽取器的输出是一个结构化对象，而不是自由文本。这个约束使后续的资格判定可以按字段检查，而不是依赖另一个LLM来判断前一个LLM的输出是否合理。解析失败的响应被记录为错误，不进入信念更新；系统不使用金标结果作为回退。这个设计是为了避免在评估中把解析失败悄悄地替换成正确答案。

锚点抽取的成本主要来自输入和输出长度。单条断言的输入通常在数千token量级，输出在数百token量级。系统按请求负载的哈希缓存响应，相同输入产生相同输出。缓存不包含API密钥，也不写入审计文件的公开部分。

\section{五类资格判定}

资格判定是FinPlanRAG与普通RAG的关键区别。普通RAG把检索到的片段直接交给生成器；FinPlanRAG在锚点进入信念更新之前，先用确定性规则检查锚点的资格。表\ref{tab:qualification}给出五类判定。

\begin{table}[H]
\centering
\zihao{5}
\caption{五类锚点资格判定。}
\label{tab:qualification}
\begin{tabular}{p{2.1cm}p{7.0cm}p{4.2cm}}
\toprule
类型 & 判定条件 & 处理 \\
\midrule
空泛 & 没有隐含标签，或锚点不提供有效标签信息 & 不进入更新 \\
过期 & 锚点日期或来源日期晚于截止时点 & 不进入更新 \\
对抗 & 来源编号不在证据集合中，或作用域不匹配 & 不进入更新 \\
硬 & 显式约束、引文通过、来源日期已知、可靠性不低于阈值 & 约束可行集 \\
软 & 其余有效锚点 & 似然更新 \\
\bottomrule
\end{tabular}
\end{table}

资格判定的具体规则如下。

第一，来源检查。锚点的来源编号必须出现在本次证据集合中。来源不存在，判为对抗。这个检查防止LLM引用一个看起来合理但实际不存在的来源。

第二，日期检查。锚点日期或来源日期晚于截止时点，判为过期。如果锚点日期与来源日期不一致，系统以来源日期为准检查过期，并对可靠性乘以0.9。日期不一致本身不是拒绝理由，但它是一个时间一致性疑点。

第三，引文检查。锚点的引文必须在来源文本中逐字出现。系统对引文和来源文本做空白规范化，然后做子串匹配。引文为空、长度过短或匹配失败，判为引文未通过。未通过的锚点可靠性乘以0.7，且不能成为硬锚点。这个检查是确定性的，不依赖另一个LLM。

第四，作用域检查。锚点的作用域必须是公司、行业、市场或监管之一。作用域不匹配，判为对抗。例如，把同行业其他公司的结果当作目标公司的直接结果，作用域不匹配。

第五，无日期检查。锚点日期和来源日期都未知时，锚点保留为软锚点，但可靠性上限为0.55，且不能成为硬锚点。这个规则承认无日期证据可能有用，但不允许它主导判断。

第六，参考值检查。对于增速类断言，只描述上一年度参考值的锚点判为空泛。上一年度实际值是契约的参考口径，不是目标期间的结果证据。这个检查防止抽取器把基数当作结果，从而把信念拉向低增长。

第七，监管规则检查。对于披露日期断言，泛泛的法定截止日判为空泛。截止日是一个上界，不能区分三个日期窗口；只有公司预约披露日期或历史披露记录才能支持具体窗口。

第八，日期标签投影。对于披露日期断言，如果锚点提到的日期与隐含标签不一致，系统把日期中的月和日投影到目标年份，再匹配日期窗口。例如，历史年报日期2025-04-21投影到2027-04-21，落在“晚于2027-04-15”窗口。投影只在一个锚点提到的所有日期指向同一窗口时生效；如果锚点同时提到多个窗口的日期，系统不覆盖模型标签，只记录不一致。这个规则修正了抽取器在跨年日期上的常见错误。

第九，时间一致性检查。对于披露日期断言，如果引文中提到的日期晚于来源发布日期，且不晚于截止时点，系统标记时间一致性疑点，可靠性乘以0.5，且不能成为硬锚点。这个检查处理动态网页或转载文章在原始发布日期之后被更新，导致引文包含当时不可能知道的信息的情况。

资格判定的输出是每条锚点的状态、资格类型、有效可靠性、引文是否通过、时间一致性标记和理由。被拒绝的锚点保留在审计文件中，理由逐条列出；被采纳的锚点进入信念更新。这个设计使审计者可以看到系统拒绝了什么，以及为什么拒绝。

\section{时间演化}

时间演化的目标是让锚点类型的可靠性随已揭示结果更新，同时不把未来信息泄漏到时点判断。设断言$q$的截止时点为$c$，锚点类型为$\tau$。截至$c$已揭示的、类型为$\tau$的锚点中，隐含标签与结果一致的数量为$m_\tau$，总数为$n_\tau$。类型$\tau$的Beta可靠性为
\begin{equation}
r_\tau=\frac{m_\tau+\alpha}{n_\tau+\alpha+\beta},
\label{eq:beta}
\end{equation}
其中$\alpha,\beta$是Beta先验参数。默认取$\alpha=\beta=1$，即均匀先验。当$n_\tau=0$时，$r_\tau=0.5$。

抽取器给出的可靠性是LLM对单条锚点的主观判断，不是经过验证的频率。系统把演化可靠性与LLM可靠性做收缩：
\begin{equation}
\rho_\tau=\lambda r_\tau+(1-\lambda)r_{\mathrm{LLM}},
\label{eq:shrink}
\end{equation}
其中$\lambda$是收缩系数，默认0.5。$\rho_\tau$作为软锚点更新的指数。这个收缩的作用是：当历史样本少时，系统更多依赖LLM判断；当历史样本多时，系统更多依赖已揭示结果的频率。

先验也按揭示日更新。设截至$c$已揭示的结果中，落在标签$h$的数量为$n_h$，总数为$n$，则先验为
\begin{equation}
b_0(h)=\frac{n_h+1}{n+|\mathcal H|}.
\label{eq:prior}
\end{equation}
当$n=0$时，先验为均匀分布。这个先验只使用已揭示结果，不使用目标期间的结果。

演化状态按轮次更新。第$r$轮的结果在第$r$轮全部断言裁决完成后才更新状态，并从第$r+1$轮开始使用。新发现的锚点类型的生效轮次为$r+1$，因此触发发现的断言不能使用该类型。这个规则保证类型库扩展是时点正确的。本文的主实验报告单轮扩展；多轮扩展作为接口保留，第\ref{ch:conclusion}章讨论其局限。

\section{类型库的验证集扩展}

固定锚点类型库依赖人工设计，难以覆盖不同指标和不同市场的证据模式。本文在Beta可靠度演化之外加入一个受验证集约束的类型库扩展步骤。它不是让模型自由改写标签或测试集，而是让模型在验证误差上提出新的锚点类型，只有通过验证集损失检验的类型才进入类型库。

设第$t$轮的锚点类型库为$\mathcal T_t$，验证集上由当前类型库得到的决策损失为$L_{\mathrm{val}}(\mathcal T_t)$。循环分四步。

第一步，误差挖掘。把验证集按公司不重叠地分成挖掘集和选择集。在挖掘集上，用当前类型库计算每条断言的预测和损失，选出预测错误或损失较高的断言。挖掘集只用于提出候选类型，不参与候选接受。

第二步，类型提出。把错误断言、对应证据、当前锚点和真实结果交给一个LLM，要求它提出至多$k$个新的锚点类型。每个候选类型是一个四元组
\begin{equation}
\tau'=(\text{名称},\ \text{L0关系},\ \text{字段},\ \text{适用条件}),
\label{eq:candidatetype}
\end{equation}
并附带一句证据示例。提出器被要求只描述可复用的证据模式，不能引用具体测试断言，不能把结果标签本身当作类型。

第三步，候选重抽取。把候选类型加入类型库，重写锚点抽取提示词中的类型说明，在选择集上重新抽取锚点并计算后验。类型库扩展为$\mathcal T_t\cup\mathcal T'$。

第四步，验证门控。只有当选择集上的损失下降超过阈值$\delta$时，才接受候选类型：
\begin{equation}
L_{\mathrm{val}}(\mathcal T_t)-L_{\mathrm{val}}(\mathcal T_t\cup\mathcal T')>\delta.
\label{eq:typeaccept}
\end{equation}
否则拒绝整组候选，保持原类型库。接受后进入下一轮，最多进行$R$轮。测试集只在循环冻结后评估一次，不参与类型提出、阈值选择或轮数选择。

这个循环与提示词演化分开评估。提示词改写、类型扩展和二者组合分别作为候选臂，在同一选择集上比较；只有损失最低且超过阈值的臂被接受。这样可以把“模型提出了新类型”和“新类型真的改善了验证损失”分开，避免把语言上的合理性当作经验收益。

类型库扩展的代价是额外的验证集抽取。每轮需要重抽选择集锚点，但不重抽测试集；类型库冻结后才抽取测试集。由于循环只使用验证集，它与主实验的时点约束一致，也不会把测试结果泄漏到类型库。

类型库扩展是框架的训练步骤，不是独立模块。算法\ref{alg:typeexpand}给出伪代码，其输出$\mathcal T^\star$作为算法\ref{alg:anchorevolve}的类型库输入。主实验中DeepSeek使用通过验证门控的类型库，GLM和Qwen保留基础库；这样，类型库的选择和其他超参数一样，只由验证集决定。

\begin{algorithm}[H]
\caption{类型库的验证集扩展}
\label{alg:typeexpand}
\begin{algorithmic}[1]
\Require 类型库$\mathcal T_0$，验证集$\mathcal D_{\mathrm{val}}$，阈值$\delta$
\Ensure 冻结类型库$\mathcal T^\star$
\State 将$\mathcal D_{\mathrm{val}}$按公司不重叠划分为挖掘集$\mathcal D_{\mathrm{mine}}$和选择集$\mathcal D_{\mathrm{select}}$
\State 在$\mathcal D_{\mathrm{mine}}$上计算当前类型库的损失，选出高损失断言
\State 由LLM提出候选类型$\mathcal T'$，每个类型包含关系、字段和适用条件
\State 在$\mathcal D_{\mathrm{select}}$上重抽锚点，计算$L_{\mathrm{val}}(\mathcal T_0\cup\mathcal T')$
\If{$L_{\mathrm{val}}(\mathcal T_0)-L_{\mathrm{val}}(\mathcal T_0\cup\mathcal T')>\delta$}
\State $\mathcal T^\star\gets\mathcal T_0\cup\mathcal T'$
\Else
\State $\mathcal T^\star\gets\mathcal T_0$
\EndIf
\State \Return $\mathcal T^\star$ \Comment{冻结后用于测试集}
\end{algorithmic}
\end{algorithm}


\section{信念更新}

信念从先验$b_0$开始。对于软锚点$a$，设其隐含标签为$\hat h$，有效可靠性为$\rho$，软似然为
\begin{equation}
L_a(h)=
\begin{cases}
p, & h=\hat h,\\
\dfrac{1-p}{|\mathcal H|-1}, & h\ne\hat h,
\end{cases}
\label{eq:softlikelihood}
\end{equation}
其中$p$是配置项，默认0.7。软锚点的更新为
\begin{equation}
b(h)\propto b(h)\,L_a(h)^{\rho}.
\label{eq:softupdate}
\end{equation}
指数$\rho$控制锚点的影响强度：可靠性越高，更新越强；可靠性越低，更新越弱。更新后对信念做归一化。

对于硬锚点，系统把信念压到隐含标签上，并保留一个小底：
\begin{equation}
b(h)=
\begin{cases}
1-\epsilon, & h=\hat h,\\
\dfrac{\epsilon}{|\mathcal H|-1}, & h\ne\hat h,
\end{cases}
\label{eq:hardupdate}
\end{equation}
其中$\epsilon$默认0.02。小底的作用是避免后续软锚点无法恢复。多个硬锚点的可行集交集、冲突检测和回退到软更新是论文的设计目标，当前实现只处理单硬锚点的点质量更新；第\ref{ch:conclusion}章把这一点列为局限。

锚点后验记为$b_A$。抽取器还输出标签分数$s$，它是边际信号，不是锚点状态的替代。系统把$b_A$与$s$按权重$\gamma$混合：
\begin{equation}
b_A'=\gamma b_A+(1-\gamma)s,
\label{eq:blend}
\end{equation}
其中$\gamma$默认0.5。混合后的后验保留锚点的结构化信息，同时利用抽取器对三类结果的边际判断。

结构化随机森林后验记为$b_R$。输入特征包括指引宽度、当前季度值相对指引中点的比例、指引与当前季度的比值、历史结果计数与频率、前次指引增长与宽度、目标季度、截止年份、行业，以及演化先验特征（同行业历史、公司意外历史、日历、当前季度相对上下界的位置、可靠性标志）。随机森林在开发集与验证集上训练，并在公司不重叠折上做温度缩放；测试公司不出现在训练集中。

锚点后验与随机森林后验融合：
\begin{equation}
b_F=w b_A'+(1-w)b_R,
\label{eq:fusion}
\end{equation}
其中$w$只在验证集上从配置的权重网格中选择，测试集不参与选择。默认网格为$\{0,0.1,0.2,0.3\}$。最终实验中选择的权重在正文中报告。

\section{全动作损失感知读出}

给定融合后验$b_F$，系统在动作空间$\mathcal A$上选择期望损失最小的动作。动作空间包括$|\mathcal H|$个点答案、$\sum_{k=2}^{|\mathcal H|}\binom{|\mathcal H|}{k}$个非单元素裁决集，以及弃权。对于三标签结果空间，动作空间包含3个点答案、4个非单元素裁决集和1个弃权，共8个动作。

点答案的期望损失为
\begin{equation}
\E[L_{\text{ans}}(h)]=\sum_{y\in\mathcal H}b_F(y)L_{\text{ans}}(h,y).
\label{eq:eans}
\end{equation}
裁决集的期望损失为
\begin{equation}
\E[L_{\text{set}}(R)]=\sum_{y\in\mathcal H}b_F(y)L_{\text{set}}(R,y).
\label{eq:eset}
\end{equation}
弃权的期望损失为$\ell_{\text{abstain}}$。系统选择
\begin{equation}
a^\star=\arg\min_{a\in\mathcal A}\E[L(a)].
\label{eq:fullaction}
\end{equation}

全动作读出与受限可信集门的区别在于动作空间。受限门只比较点答案和最高密度$1-\alpha$可信集，可信集的大小由覆盖率目标决定，而不是由损失决定。全动作读出把点答案、裁决集和弃权放在同一个损失下比较，因此可以在不确定性高时选择裁决集或弃权，也可以在置信度高时选择点答案。对于同一个后验和同一个损失，全动作读出的期望损失不会高于受限门，因为受限门的动作是全动作空间的子集。这个性质是H3的形式基础。

单元素裁决集不进入动作空间。原因是单元素裁决集与点答案在语义上重复，但在损失上不同：如果漏报惩罚低于错误分类惩罚，单元素裁决集会系统性地优于点答案，导致系统把点答案都标成裁决集。这个现象不是决策改进，而是动作标签的重复计数。本文把裁决集限制为两个及以上标签，使三类动作的含义保持清晰。

可信集$R_\alpha$仍然作为覆盖口径报告：
\begin{equation}
R_\alpha=\arg\min_{R\subseteq\mathcal H}\{|R|:\sum_{h\in R}b_F(h)\ge1-\alpha\}.
\label{eq:credible}
\end{equation}
它不参与动作选择，但用于审计和覆盖率报告。损失最优裁决集$R^\star$是动作选择的结果，可以在$R_\alpha$更宽时选择更窄的集合，也可以在不确定性高时选择弃权。两者在审计报告中分开列出，避免把覆盖率误读为决策质量。

\section{披露日期断言}

披露日期断言的算法与数值指标断言共享同一套锚点状态和读出，但结果空间和日期估计不同。

结果空间由日历窗口定义。给定目标期间$p$，系统根据交易所规则和公司历史披露习惯构造三个窗口。例如，2026年年度报告的窗口为“不晚于2027-03-31”“2027-04-01至2027-04-15”“晚于2027-04-15”。窗口必须互斥且覆盖法定披露截止日之前的全部日期。

锚点抽取器除了输出锚点和标签分数，还输出一个日期估计，包括点估计、区间下界、区间上界和依据。日期估计是辅助信息，不直接决定动作；动作仍由标签后验和损失读出决定。系统把日期估计与标签后验做一致性检查：如果点估计落在最高信念窗口内，系统保留点估计；如果点估计落在窗口外，系统用最高信念窗口的中点作为回退点估计。区间下界和上界优先使用模型给出的值，但限制在合理范围内。

日期标签投影是披露日期断言的关键修正。历史年报日期、预约披露日期和监管截止日往往跨年，抽取器容易把历史日期直接映射到当前窗口。系统把日期中的月和日投影到目标年份，再匹配窗口。投影规则只在一个锚点提到的所有日期指向同一窗口时生效，避免把包含多个日期的锚点强行归类。

披露日期断言的评价不以点日期误差作为主指标，因为公告日期与实际披露时点可能相差一天。本文在A股回测上报告日期窗口准确率，并在案例研究中展示动作、裁决集、点估计、区间和审计轨迹。

\section{算法与复杂度}

算法\ref{alg:anchorevolve}给出单条断言的完整流程。

\begin{algorithm}[H]
\caption{AnchorEvolve单条断言的证据构造、锚点演化与全动作裁决}
\label{alg:anchorevolve}
\begin{algorithmic}[1]
\Require 断言契约$q=(e,p,m,\mathcal H,c)$，演化状态$S_c$，类型库$\mathcal T^\star$，随机森林后验$b_R$
\Ensure 动作$a^\star$，点答案$h^\star$，裁决集$R^\star$，后验$b_F$
\State $V_c\gets\textsc{CollectEvidence}(q,c)$ \Comment{按截止时点过滤，分配$W\#$}
\State $A\gets\textsc{ExtractAnchors}(q,V_c,S_c,\mathcal T^\star)$ \Comment{LLM，要求来源与逐字引文}
\State $A_q\gets\textsc{Qualify}(A,q,c)$ \Comment{空泛/过期/对抗/硬/软}
\State $b_A\gets\textsc{BeliefUpdate}(S_c,A_q)$
\State $b_A'\gets\gamma b_A+(1-\gamma)\textsc{LabelScores}(A)$
\State $b_F\gets w b_A'+(1-w)b_R$
\State $(a^\star,h^\star,R^\star)\gets\textsc{FullActionReadout}(b_F)$
\State $S_{c^+}\gets\textsc{RevealUpdate}(S_c,c)$ \Comment{仅用揭示日早于下一截止时点的结果}
\State \Return $(a^\star,h^\star,R^\star,b_F)$
\end{algorithmic}
\end{algorithm}

算法的时间成本由三部分组成。第一是证据构造，包括网页检索、页面抓取和官方公告查询。网页检索的数量由配置限制，默认每条断言四条查询；页面抓取只抓前若干条结果，避免不必要的网络开销。第二是锚点抽取，一次LLM调用的输入在数千token量级，输出在数百token量级。第三是读出，动作空间的大小为$|\mathcal H|+\sum_{k=2}^{|\mathcal H|}\binom{|\mathcal H|}{k}+1$，当$|\mathcal H|\le6$时最多63个动作，计算可以忽略。随机森林推断在本地完成，不产生API成本。

系统的存储成本主要来自审计文件。每条断言保存证据列表、锚点列表、信念轨迹、动作损失和LLM响应摘要。审计文件的大小通常在数百KB量级，适合长期保存和人工复核。LLM响应缓存单独存放，不进入公开审计文件。

\section{可审计性}

可审计性是FinPlanRAG的设计目标之一。系统把每个中间状态写入结构化文件，使读者可以从最终动作回溯到证据和锚点。表\ref{tab:auditfields}给出审计文件的主要字段。

\begin{table}[H]
\centering
\zihao{5}
\caption{审计文件的主要字段。}
\label{tab:auditfields}
\begin{tabular}{p{2.8cm}p{8.8cm}}
\toprule
对象 & 字段 \\
\midrule
契约 & 公司、代码、目标期间、指标、结果空间、标签定义、参考口径、截止时点 \\
证据 & 来源编号、标题、链接、发布日期、提供者、摘要、正文、是否官方来源 \\
锚点 & 类型、取值、来源编号、逐字引文、发布日期、可靠性、资格、隐含标签、理由、纠正记录 \\
信念 & 先验、仅锚点后验、融合后验、逐步更新轨迹 \\
动作 & 动作类型、点答案、裁决集、可信集、期望损失、各动作损失排序 \\
日期 & 点估计、区间、最高信念窗口、模型估计与窗口是否一致、依据 \\
\bottomrule
\end{tabular}
\end{table}

审计文件有两个用途。第一，复核。分析师可以逐条检查被采纳的锚点是否真的支持隐含标签，被拒绝的锚点是否应当被拒绝，证据的发布日期是否早于截止时点。第二，问责。当判断出现错误时，系统可以定位错误来自证据缺失、锚点抽取、资格判定、信念更新还是动作选择，而不是把错误归因于一个黑箱模型。

可审计性不等于正确性。一个锚点即使有来源和引文，也可能是对来源的误读；一个资格判定即使规则明确，也可能过严或过松。本文在案例研究中展示审计轨迹，并在第\ref{ch:conclusion}章把人工审计实验列为未来工作。当前的可审计性是设计属性，不是已经测量的用户收益。
